\documentclass[10pt,a4paper]{amsart}
\usepackage[margin=19mm]{geometry}
\usepackage{amsmath,amssymb,mathtools}
\usepackage[scr=rsfs,cal=euler]{mathalfa}
\usepackage{xfrac}
\usepackage[unicode,hidelinks,bookmarks=false]{hyperref}
\newcommand{\ie}{i.e.\ }
\newcommand{\cf}{cf.\ }
\newcommand{\resp}{respectively\ }
\newcommand{\Subsection}{\subsection}
\providecommand{\nobreakdash}{\nobreak\hbox{-}}
\setlength{\emergencystretch}{2em}
\multlinegap0pt
\usepackage{amssymb}
\usepackage[shortlabels]{enumitem}
\newtheorem{lemma}{Lemma}
\title{\bf On irreducible factors of polynomials over integers}
\author{Rishu Garg, Jitender Singh}
\date{Source: arXiv:2512.20262v1 (CC BY 4.0)}
\begin{document}
\maketitle

\noindent\emph{Source and licence.} \bf On irreducible factors of polynomials over integers. Version arXiv:2512.20262v1 (CC BY 4.0), distributed by arXiv under the Creative Commons Attribution 4.0 International licence, http://creativecommons.org/licenses/by/4.0/, as recorded in the arXiv licence metadata for this exact version and shown on the abstract page https://arxiv.org/abs/2512.20262v1. Full licence notice: \texttt{licenses/arxiv-2512.20262-CC-BY-4.0.txt}.\par\smallskip

\noindent\textit{Editorial scope.} Counted unit: the article's lemma L:5 (source0.tex lines 448--454). The article's complete original proof of it is delivered with it (source0.tex lines 468--494, one \texttt{proof} environment of 27 lines, which the article places away from the statement). No mathematics is written, repaired or re-derived: the statement and the proof are the article's own lines, in the article's own notation, and no retained line is substituted or re-flowed.  No further source-local context is needed: every cross-reference of every retained line resolves inside the two delivered spans.  Nothing was omitted from any retained span, so the package carries no omission record.  No retained line contains a citation, so the wrapper carries no bibliography and no bibliography entry.  No retained line contains a figure environment, an image-inclusion command or drawing code, so no image is delivered and the package declares no asset dependency.  Editorial formatting only, in addition: an AMS \texttt{amsart} preamble that restates the article's own declarations and macros used by the retained lines, namely the environments \texttt{lemma} on separate counters, together with the standard packages those lines need.  The retained environments print in this document as Lemma 1 in order of appearance, whereas the article prints the same items with the numbering of its own full document.  The complete native source of the article as distributed in the arXiv e-print for this exact version is bundled unchanged as \texttt{sources/191548/source0.tex}.
\begin{lemma}\label{L:5}
Let $g = s_0 + s_1z +\cdots+s_n z^n\in \Bbb{Z}[z]$ be a primitive polynomial with $s_0s_n\neq 0$ such that $s_n=\pm p^k d$ for some  positive integers $k$ and $d$ and a prime $p\nmid d$. Suppose that all zeros of $g$ lie in the region $|z|>\sqrt[\Delta_f]{d}$ in the complex plane.  If $|s_0/q|\leq |s_n|$, where $q$ is the smallest prime divisor of $s_0$, then the polynomial $g$ is a product of at most
\begin{eqnarray*}
\min_{1\leq i\leq n}\{k,i+v_p(s_{n-i})\}
\end{eqnarray*}
irreducible factors in $\mathbb{Z}[z]$. In particular, if $k=1$, or $p\nmid s_{n-1}$, then $g$ is irreducible.
\end{lemma}

\begin{proof}[\bf Proof of Lemma \ref{L:5}]
Let $g(z)=g_1(z)\cdots g_N(z)$ be a product of $N$ irreducible polynomials $g_1,\ldots,g_N$ in $\mathbb{Z}[z]$. Then $\Delta_f\leq \deg g_i$ for all $i=1,\ldots,N$. Since $1\leq \min_{1\leq i\leq n}\{i+v_p(s_{n-i})\}$, we may assume without loss of generality that $N>1$. Let $\alpha_i\neq 0$ be the leading coefficient of $g_i$ for each $i$. We then have from the hypothesis that $\pm p^kd=s_n=\alpha_1\cdots \alpha_N$. We may write $g_i(z)=\alpha_i\prod_{\theta}(z-\theta)$ where the product runs over all zeros $\theta$ of $g_i$ for each $i$. We then arrive at the following:
\begin{eqnarray*}
|g_i(0)| &=& |\alpha_i|\prod_{\theta}|\theta|>  |\alpha_i|\bigl(\sqrt[\Delta_f]{d}\bigr)^{\deg g_i}\geq |\alpha_i|\bigl(\sqrt[\Delta_f]{d}\bigr)^{\Delta_f}=|\alpha_i|d\geq d,
\end{eqnarray*}
which shows that $|g_i(0)|>d$, and so, $|g_i(0)/q|\geq 1$ for all $i=1,\ldots, N$. Consequently, using the hypothesis, we get
\begin{eqnarray*}
|s_0/q| \leq |s_n|&<&|s_n|\left({\frac{|g_1(0)|}{|\alpha_1|d^{\frac{\deg g_1}{\Delta_f}}}\cdots \frac{|g_N(0)|}{|\alpha_N|d^{\frac{\deg g_N}{\Delta_f}}}}\right)\times \frac{|\alpha_i|d^{\frac{\deg g_i}{\Delta_f}}}{|g_i(0)|}
= |s_0|\frac{|\alpha_i|}{|g_i(0)|d^{\frac{n-\deg g_i}{\Delta_f}}},
\end{eqnarray*}
for each fixed index $i$, which yields $|\alpha_i|> |g_i(0)/q|d^{\frac{n-\deg g_i}{\Delta_f}}$ for each such $i$. We also
observe that
\begin{eqnarray*}
n-\deg g_i=\sum_{j=1,j\neq i}^{N} \deg g_j\geq \sum_{j=1,j\neq i}^{N}\Delta_f=(N-1)\Delta_f,
\end{eqnarray*}
for each index $i$, which in view of the fact that $N>1$ tells us that $(n-\deg g_i)/\Delta_f\geq N-1\geq 1$. This observation along with that $|g_i(0)/q|\geq 1$, we have
\begin{eqnarray*}
|\alpha_i|>|g_i(0)/q|d^{\frac{n-\deg g_i}{\Delta_f}}\geq d^{\frac{n-\deg g_i}{\Delta_f}}\geq d,~i=1,\ldots,N.
\end{eqnarray*}
Consequently, we have $|\alpha_i|>d$ for each $i$. This in the view that $p^k d=|s_n|=|\alpha_1|\cdots |\alpha_N|$ and the hypothesis that $p\nmid d$ shows that $p$ divides $|\alpha_i|$ for each $i=1,\ldots, N$. This proves that $N\leq k$.

Now let $j$ be a positive integer with $j\leq n$, and let $v_p(s_{n-j})=J\geq 0$. We may assume that $k > j+J$. Assume on the contrary that $N>j+J$, that is $N-j> J\geq 0$.  We may express $g_i=\sum_{t=0}^{\deg g_i}b_{it}z^t$ so that $\alpha_i=b_{i\deg g_i}$ for each $i$. Since we have $g(z)=g_1(z)\cdots g_N(z)$, we have $n=\sum_{i=1}^N \deg g_i$. Consequently, we have
\begin{eqnarray*}
s_{n-j}&=&\sum_{i_1+\cdots+i_N=n-j}b_{1i_1}\cdots b_{Ni_N}=\sum_{\sum_{t=1}^N(\deg g_t-i_t)=j} b_{1i_1}\cdots b_{Ni_N},
\end{eqnarray*}
where $0\leq i_t\leq \deg g_t$ and $0\leq \deg g_t -i_t\leq j$ for each index $i_t$. Since $N>j+J$, it follows that at least $N-j$ of the numbers $\deg g_1-i_1,\ldots,\deg g_N-i_N$ must be each equal to zero for any choice of the $N$-tuple $i_1,\ldots,i_N$. In view of this, for any $N$-tuple $i_1,\ldots, i_N$, at least $N-j$ of the indices satisfy $i_t=\deg g_t$, $1\leq t\leq N$, so that $b_{ti_t}=b_{t\deg g_t}=\alpha_t$, the leading coefficient of $g_t$, where we know that $p\mid \alpha_t$. Thus $p^{N-j}$ divides each term of the form $b_{1i_1}\cdots b_{Ni_N}$, and so, $p^{N-j}$ divides $\displaystyle \sum_{i_1+\cdots+i_N=n-j}b_{1i_1}\cdots b_{Ni_N}=s_{n-j}$, and so, $J=v_p(s_{n-j})\geq N-j>J$, which is a contradiction.
\end{proof}

\end{document}
